\section{Theory: modular addition}
\subsection{Setting}
\label{sec:setting}

\paragraph{Modular Addition} In modular addition for an odd prime $p$ with $n\ge 1$ operands, the input is an
$n$-tuple $a\in\Zp^n$ of residues mod $p$; the result is $s(a)=\sum_{i=1}^n a_i \bmod p$.


\paragraph{Margin} On input $a$ the model outputs one score per candidate answer
$c\in\Zp$; write $Z(a,c)$. So $Z$ is a table with one row per input and one column per candidate.
The \emph{margin} is the worst-case gap between the right answer and the best wrong one,
\begin{equation}
  \marg(Z) \;=\; \min_{a}\ \min_{c\neq s(a)}\ \big[Z(a,s(a)) - Z(a,c)\big].
\end{equation}

\paragraph{Neuron-additive homogeneous models.} We consider models in which the table is a sum of
per-neuron tables, each neuron paying its own cost:
\begin{equation}
  z=\sum_{i=1}^m z_i,\qquad \Pq(\theta)=\sum_{i=1}^m \|\theta_i\|^q ,
\end{equation}
where $\theta_i$ is the weight vector of neuron $i$, $z_i$ depends on $\theta_i$ only, and
$z_i$ is homogeneous of degree $\nu$: scaling $\theta_i$ by $t$ scales $z_i$ by $t^\nu$. In a
one-hidden-layer network $z_i$ is the table produced by neuron $i$ alone; the property fails when a
shared layer precedes the neurons. The exponent $q\in[2,\nu]$ sets the penalty: $q=2$ is plain
weight decay, $\lambda\|\theta\|^2$; $q=\nu$ is the $L_{2,\nu}$ penalty of \citet{morwani2024feature}.
Two derived numbers run the section,
\begin{equation}
  r:=q/\nu\in(0,1],\qquad s:=\nu/q=1/r\ge 1 :
\end{equation}
$r$ is how much cheaper a feature gets per unit of amplitude as it grows, and $s$ is the return on
concentrating budget into few features ($s=1$, Morwani's penalty, means no return).

\paragraph{Normalised margin.} $\gam(\theta)=\marg(z)/\Pq(\theta)^{\nu/q}$ is unchanged when all
weights are scaled, so it depends only on the weight \emph{direction}. Its largest value over all
directions is the quantity the section characterises. By \citet[Thm.~4.1]{wei2019regularization}, the
global minimiser of cross-entropy plus $\lambda\Pq$ approaches a direction attaining that largest
value as $\lambda\to 0$; nothing here concerns whether gradient descent finds it.

\paragraph{Instances.} The \emph{product model} $z_i(a,c)=w_i(c)\prod_j e_{ij}(a_j)$ (one lookup
vector per operand plus a readout; $\nu=n+1$); the \emph{power model}
$z_i(a,c)=w_i(c)\big(\sum_j e_{ij}(a_j)\big)^n$ of \citet{morwani2024feature,li2025nsum} ($\nu=n+1$);
and the one-hidden-layer \emph{ReLU model} ($\nu=2$). Appendix~\ref{app:models} checks the
assumptions for each.

\subsection{Clean part, features, and the programme}
\label{sec:objects}

The score of the candidate sitting $D$ steps below the right answer, averaged over all inputs, is
\begin{equation}
  \clean(D)\;=\;\frac{1}{p^n}\sum_{a} z\big(a,\ s(a)-D\big),\qquad D\in\Zp .
\end{equation}
$\clean$ is the part of the table that depends only on (answer $-$ candidate). Everything in $z$
that is not of this form is \emph{interference}: content that does not depend on whether the
candidate is right. Since $\clean$ is a function on the cyclic group $\Zp$ it is a sum of cosines,
\begin{equation}
  \clean(D)\;=\;\bar{\clean}+\sum_{k\in\Fp} A_k\cos\!\big(2\pi kD/p+\phi_k\big),\qquad
  \Fp=\{1,\dots,\tfrac{p-1}{2}\},\ A_k\ge 0,
  \label{eq:spectrum}
\end{equation}
and we call $A_k$ the \emph{clean amplitude} and $\phi_k$ the \emph{phase} of frequency $k$. A
\emph{feature} is a frequency with $A_k>0$; the \emph{feature count} of a table is the number of its
features. This is the sense in which \citet{nanda2023progress} and \citet{morwani2024feature} count
features. Because $z=\sum_i z_i$ and $\clean$ is linear in $z$, each neuron has its own amplitudes and
phases $(a_{ik},\phi_{ik})$ and
\begin{equation}
  A_k e^{i\phi_k}=\sum_i a_{ik}e^{i\phi_{ik}},\qquad\text{so}\qquad A_k\le\sum_i a_{ik}.
  \label{eq:atoms}
\end{equation}
We call $c_{ik}:=a_{ik}e^{i\phi_{ik}}$ the \emph{atom} of neuron $i$ at frequency $k$.

Frequency $k$ separates the right answer from the wrong answer $D$ steps away by the
\emph{coverage}
\begin{equation}
  \cov{k}{D}=1-\cos(2\pi kD/p)\ \ge 0,\qquad \sum_{D\neq 0}\cov{k}{D}=p .
\end{equation}
A single cosine barely separates offsets with $kD\approx 0 \bmod p$; several frequencies are needed
so that every offset is covered by someone. The \emph{programme} is
\begin{equation}
  \Gam_s(p)\;=\;\max\Big\{\ \min_{D\neq 0}\sum_{k\in\Fp} y_k\,\cov{k}{D}\ :\ y\ge 0,\ \sum_{k\in\Fp} y_k^{1/s}\le 1\ \Big\}.
  \label{eq:programme}
\end{equation}
Read it as: split one unit of budget over frequencies, where a budget share $y_k^{1/s}$ buys clean
amplitude $y_k$, collect the worst-case coverage over offsets, and maximise. The maximum exists (compact
constraint set, continuous objective). We write $\Kstar(s,p)=|\supp(y^\star)|$ for the support size of a
maximiser $y^\star$: the \emph{feature count of the programme}. Restricting to equal shares on a $K$-set,
with $M_K(p)=\max_{|S|=K}\min_{D\neq 0}\sum_{k\in S}\cov{k}{D}$, gives the lower bound
$\Gam_s(p)\ge\max_K K^{-s}M_K(p)$; the ladder $K^{-s}M_K$ is a cheap approximation to the programme.

\subsection{The ansatz, derived lightly}
\label{sec:ansatz}

A feature of a degree-$\nu$ model is carried by $\nu$ weight vectors along the path from input to
logit. If they have norms $\rho_1,\dots,\rho_\nu$ the amplitude is proportional to $\prod_l\rho_l$ and
the penalty is $\sum_l\rho_l^q$; at fixed product the sum is smallest when all norms are equal
(AM--GM), so the cost of amplitude $A$ is proportional to $A^{q/\nu}=A^{1/s}$. Because $q\le\nu$ the
exponent is at most one: amplitude is cheaper per unit the more of it one feature carries. The
margin is worst-case coverage, $\min_D\sum_k y_k\cov{k}{D}$. Maximising margin per unit cost is
therefore exactly \eqref{eq:programme}, and with equal amplitudes on a $K$-set it is $M_K/K^s$:
coverage grows with $K$, the concentration price $K^{-s}$ falls with $K$, and the feature count is
where the product peaks. At $s=1$ the budget is linear and the maximiser is every frequency at equal
share (Corollary~\ref{cor:dense}); for $s>1$ the dense solution loses, since $M_K\le K$ gives
$K^{-s}M_K\to 0$, and the argmax is a finite $\Kstar$ that falls as $s$ grows, i.e.\ as the degree
rises or the penalty exponent falls. What the heuristic hides is (i) that the margin can be bounded by
the clean amplitudes alone although trained tables carry interference, and (ii) that one neuron does not
produce a clean feature but a feature plus sign-flipped copies of it, whose cancellation costs a
factor at $q<\nu$. The rest of the section supplies both.

\subsection{Three lemmas for any table}
\label{sec:lemmas}

\begin{lemma}[interference never helps]
\label{lem:junk}
For every table $z$ on the full table and every $D\neq 0$,
$\ \marg(z)\le\clean(0)-\clean(D)$.
\end{lemma}
\begin{proof}
For every input $a$ the candidate $s(a)-D$ is wrong, so $z(a,s(a))-z(a,s(a)-D)\ge\marg(z)$. Average
over all $p^n$ inputs.
\end{proof}
The worst input is no better than the average input at any fixed offset, and only the clean part
survives the average. The proof uses every input, which is why all statements are full-table
statements.

\begin{lemma}[phase lemma]
\label{lem:phase}
For every table $z$ and every $D\neq 0$,
\begin{equation}
  \marg(z)\;\le\;\sum_{k\in\Fp} A_k\,\cov{k}{D},
  \qquad\text{hence}\qquad
  \marg(z)\;\le\;\min_{D\neq 0}\sum_{k\in\Fp} A_k\,\cov{k}{D}.
\end{equation}
\end{lemma}
\begin{proof}
By \eqref{eq:spectrum}, $\clean(0)-\clean(D)=\sum_k A_k[\cos\phi_k-\cos(2\pi kD/p+\phi_k)]$. This is not
bounded by $\sum_k A_k\cov{k}{D}$ for a single $D$, but the average of the bracket at $D$ and $-D$ is
$\cos\phi_k\,\cov{k}{D}\le\cov{k}{D}$. Apply Lemma~\ref{lem:junk} at $D$ and at $-D$ and average.
\end{proof}
Equality holds iff every $\phi_k=0$: a network gains nothing from misaligned phases, and the
bound depends on the network only through its amplitude vector $(A_k)$.

\begin{proposition}[coverage bound]
\label{prop:coverage}
With $\|A\|_r:=\big(\sum_{k\in\Fp}A_k^r\big)^{1/r}$ and $s=1/r$,
$\ \marg(z)\le\Gam_s(p)\,\|A\|_r$.
\end{proposition}
\begin{proof}
If $\|A\|_r=0$ use Lemma~\ref{lem:phase}. Otherwise $y_k=A_k/\|A\|_r$ is feasible for
\eqref{eq:programme} ($\sum_k y_k^{1/s}=\sum_k y_k^{r}=1$), so $\min_D\sum_k y_k\cov{k}{D}\le\Gam_s(p)$;
multiply by $\|A\|_r$.
\end{proof}

\subsection{The universal bound and the dense endpoint}
\label{sec:universal}

Nothing about the model has been used so far. The one thing a model must supply is how much clean
amplitude a neuron can buy.

\begin{assumption}[budget]
\label{ass:budget}
There is a constant $C_q>0$ such that for every neuron $i$,
$\ \sum_{k\in\Fp}a_{ik}^{\,q/\nu}\le C_q\,\|\theta_i\|^q$.
\end{assumption}
One neuron of weight norm $\|\theta_i\|$ can create only so much clean amplitude, priced with exponent
$q/\nu$, however it spreads it over frequencies. For the product and power models
$C_q^{\nu/q}=B:=2(2p\nu)^{-\nu/2}$ up to a model constant (Appendix~\ref{app:models}).

\begin{theorem}[universal margin bound]
\label{thm:universal}
For every neuron-additive $\nu$-homogeneous model satisfying Assumption~\ref{ass:budget}, every
width, and every $q\in[2,\nu]$,
\begin{equation}
  \gam(\theta)\;\le\;C_q^{\,\nu/q}\ \Gam_{\nu/q}(p).
\end{equation}
\end{theorem}
\begin{proof}
Let $r=q/\nu\le 1$. Since $t\mapsto t^r$ is subadditive and $A_k\le\sum_i a_{ik}$ by \eqref{eq:atoms},
$\sum_k A_k^r\le\sum_i\sum_k a_{ik}^r\le C_q\Pq(\theta)$ by Assumption~\ref{ass:budget}. Raise to the
power $1/r=\nu/q$, insert into Proposition~\ref{prop:coverage}, and divide by $\Pq^{\nu/q}$.
\end{proof}
No hypothesis beyond the budget is used: neurons may mix frequencies, carry interference, or be
inactive.

\begin{corollary}[the dense endpoint]
\label{cor:dense}
$\Gam_1(p)=p/(p-1)$, with the unique maximiser $y_k=2/(p-1)$ for all $k\in\Fp$. Hence at $q=\nu$ any table
attaining Theorem~\ref{thm:universal} has all $A_k$ equal and all phases aligned; for the product model
the bound is attained (Proposition~\ref{prop:block} with $r=1$), so the $L_{2,\nu}$ maximum-margin
product network uses every frequency.
\end{corollary}
\begin{proof}
For feasible $y$, the minimum over $D$ is at most the average over $D\neq 0$, which is
$\frac{p}{p-1}\sum_k y_k\le\frac{p}{p-1}$. Equality forces $\sum_k y_k=1$ and
$\sum_k y_k\cos(2\pi kD/p)$ constant on $D\neq 0$; a function equal to $\kappa$ off zero and $1$ at zero
has all non-zero Fourier coefficients equal, so $y$ is flat (Appendix~\ref{app:dense}).
\end{proof}
This is the theorem of \citet{morwani2024feature} recovered as the $s=1$ endpoint of the programme. For
$q<\nu$ Theorem~\ref{thm:universal} still holds but is \emph{not} attained: a single neuron per frequency
carries interference (\S\ref{sec:class}), so its atoms cannot be clean and aligned at once.
Theorem~\ref{thm:ceiling} gives the exact value on the class where interference is absent.

\subsection{The interference-free class}
\label{sec:class}

Fix a frequency $k$. A neuron carrying the cosine of $k(a_1+\dots+a_n-c)$ in general also carries
\emph{sign-flipped} copies: for a non-empty set $T$ of operands, the cosine of
$k\big(\sum_{j\notin T}a_j-\sum_{j\in T}a_j-c\big)$. Such a mode depends on inputs and candidate but not
on whether the candidate is right; it is interference in the sense of Lemma~\ref{lem:junk}, and it sits
at the same frequency as the feature. Write $J_{k,T}$ for the table's (complex) coefficient at that mode,
normalised like the atoms, and $J_{ik,T}$ for neuron $i$'s; then $J_{k,\emptyset}=A_ke^{i\phi_k}$ and
$J_{k,T}=\sum_iJ_{ik,T}$.

\begin{assumption}[sign rule]
\label{ass:sign}
For every neuron $i$ and frequency $k$ there are unimodular $v_{ijk}$ ($j=1..n$) with
$\ J_{ik,T}=c_{ik}\prod_{j\in T}v_{ijk}$ for every $T\subseteq\{1..n\}$.
\end{assumption}
A neuron's sign-flipped copy of its feature has the same size as the feature and a phase obtained by
flipping the phases of the operands in $T$. This holds whenever the neuron's contribution at these modes
is a product of one function per operand and one of the candidate (Appendix~\ref{app:models}).

\begin{definition}[interference-free class]
A table is \emph{interference-free} if $J_{k,T}=0$ for every $k\in\Fp$ and every non-empty $T$.
\end{definition}
Equivalently, at every frequency the table's content at the $2^\nu$ modes with entries $\pm k$ is a
function of (answer $-$ candidate) alone, i.e.\ it is invariant under the task's symmetry
$(a,c)\mapsto(a+t,\,c+\sum_jt_j)$. The class is a property of the \emph{table}: its neurons may mix
frequencies. Interference at other modes is not restricted; Lemma~\ref{lem:junk} already discounts it
and the budget already prices it.

Under the sign rule, an interference-free feature at $k$ is a finite system that involves neither $p$
nor the task: relabel the neurons with $c_{ik}\neq 0$ as $l=1..N$ and write $A=\sum_lc_l$; then
\begin{equation}
  \sum_{l=1}^N c_l\prod_{j\in T}v_{lj}\;=\;A\,[T=\emptyset]\qquad\text{for every }T\subseteq\{1..n\}.
  \label{eq:JF}
\end{equation}
$N$ complex numbers with phases attached, whose plain sum is $A$ and whose sum twisted by any non-empty
pattern of flips is zero.

\begin{theorem}[feature cost]
\label{thm:cost}
Let $v_{lj}$ have modulus $1$ and let $c_1,\dots,c_N$ satisfy \eqref{eq:JF} with total $A$. Then for every
$r\in(0,1]$,
\begin{equation}
  \sum_{l=1}^N|c_l|^r\;\ge\;\nu^{1-r}|A|^r,\qquad\nu=n+1 .
\end{equation}
For $r\in(0,1)$ equality holds iff exactly $\nu$ of the $c_l$ are non-zero, all of modulus $|A|/\nu$,
with pairwise orthogonal chain vectors $m_l=(1,v_{l1},v_{l1}v_{l2},\dots,v_{l1}\cdots v_{ln})$.
\end{theorem}
To carry clean amplitude $A$ while cancelling every sign-flipped copy, a feature must be spread over at
least $\nu$ neurons, evenly, and extra neurons never make it cheaper. At $r=1$ the inequality is the
triangle inequality and every aligned configuration attains it, which is why the neuron count per
feature is pinned only for $q<\nu$.

\begin{proof}[Proof sketch; full proof in Appendix~\ref{app:cost}]
Map each neuron to its chain vector $m_l\in\mathbb{C}^\nu$. The matrix $M=\sum_lc_lm_lm_l^*$ has
entry $(k,k')$, $k>k'$, equal to the sum in \eqref{eq:JF} for the interval $T=\{k'+1..k\}$, hence
zero, and diagonal $A$: it is triangular with all eigenvalues $A$. Expanding $\det(I+tM)$ by
principal minors (Cauchy--Binet) and bounding Gram determinants of unit vectors by one (Hadamard) gives
$e_k(x)\ge\binom{\nu}{k}$ for the elementary symmetric polynomials of $x_l=\nu|c_l|/|A|$, hence
$\prod_l(1+x_lt)\ge(1+t)^\nu$. Integrating the logarithm against $t^{-1-r}dt$ and using
$x^r=\frac{r\sin\pi r}{\pi}\int_0^\infty\log(1+xt)\,t^{-1-r}dt$ yields $\sum_lx_l^r\ge\nu$.
\end{proof}

\begin{theorem}[margin ceiling in the interference-free class]
\label{thm:ceiling}
For every neuron-additive $\nu$-homogeneous model satisfying Assumptions~\ref{ass:budget}
and~\ref{ass:sign}, every width, every $q\in[2,\nu]$, and every interference-free table,
\begin{equation}
  \gam(\theta)\;\le\;\nu^{-(\nu-q)/q}\ C_q^{\,\nu/q}\ \Gam_{\nu/q}(p).
\end{equation}
For the product model the bound is attained (Proposition~\ref{prop:block}).
\end{theorem}
\begin{proof}
Let $r=q/\nu$. At each $k$ Theorem~\ref{thm:cost} gives $\sum_ia_{ik}^r\ge\nu^{1-r}A_k^r$; sum over $k$
and apply the budget neuron by neuron: $\nu^{1-r}\sum_kA_k^r\le C_q\Pq$. Insert into
Proposition~\ref{prop:coverage}; $(r-1)/r=-(\nu-q)/q$.
\end{proof}
The class bound differs from the universal bound by exactly $\nu^{(\nu-q)/q}$ ($\sqrt3$ at $\nu=3$, $4$
at $\nu=4$, $5^{3/2}$ at $\nu=5$ under plain weight decay, $1$ at $q=\nu$): the price of splitting each
feature over $\nu$ neurons to cancel its interference.

\begin{corollary}[feature count and structure]
\label{cor:count}
Let $q<\nu$. Every interference-free table attaining Theorem~\ref{thm:ceiling} has clean amplitudes
$A_k\propto y^\star_k$ for a maximiser $y^\star$ of $\Gam_{\nu/q}(p)$, hence exactly $\Kstar(\nu/q,p)$
features; at each feature exactly $\nu$ neurons carry a non-zero atom, all of modulus $A_k/\nu$, with
orthonormal chain vectors.
\end{corollary}

\begin{proposition}[the $\nu$-block attains it]
\label{prop:block}
For the product model, fix $k\in\Fp$, $\rho>0$, and for $l=0..n$ let $\varphi_l=\pi l/\nu$,
$\theta_j=2\pi ka_j/p$, $\theta_c=2\pi kc/p$. The $\nu$ neurons
$e_{lj}(a)=\rho\cos(\theta_j+\varphi_l)$ (all $j$), $w_l(c)=\rho\cos(\theta_c+n\varphi_l)$ have total table
$\nu\rho^\nu2^{-n}\cos(\theta_1+\dots+\theta_n-\theta_c)$, a function of (answer $-$ candidate) alone,
and each costs $\|\theta_l\|^2=\nu\rho^2p/2$. If $y^\star$ maximises $\Gam_{\nu/q}(p)$, placing at each
$k\in\supp(y^\star)$ a block with $\rho_k^\nu=2^ny^\star_k/\nu$ gives an interference-free network with
$\gam=\nu^{-(\nu-q)/q}\,B\,\Gam_{\nu/q}(p)$.
\end{proposition}
The phases $\varphi_l$ are a rotated copy of the $\nu$-th roots of unity: the equality case of
Theorem~\ref{thm:cost}. The share weighting $\rho_k^\nu\propto y^\star_k$ is essential; the optimal
shares are not equal (at $p=59$, $s=3/2$ they range over $0.15$--$0.25$), and an equal-amplitude block on
the optimal support does not attain the bound. Proof in Appendix~\ref{app:block}.

\paragraph{What the theorems say together.} The programme and its support size $\Kstar$ are the same
objects in Theorems~\ref{thm:universal} and~\ref{thm:ceiling}; only the constant differs. Under plain
weight decay the count of the class optimum falls with the degree and grows only slowly with the
modulus (Table~\ref{tab:kstar}). The dense solution of Corollary~\ref{cor:dense}, evaluated by the same
machinery, reaches $0.52$--$0.63$ of the class optimum at $\nu=3$, $0.15$--$0.25$ at $\nu=4$ and
$0.04$--$0.08$ at $\nu=5$: the preference for sparsity sharpens with every rung of the degree.

\begin{table}[t]
\centering
\caption{$\Kstar(\nu/q,p)$ at plain weight decay ($q=2$), from the free-share programme. Support
enumeration is exhaustive for $K\le 8,7,5$ at $p=59,71,113$ and heuristic above, so a missed support can
only raise $\Kstar$. The equal-share ladder argmax agrees except at $(59,3)$: 6 and $(113,3)$: 5.}
\label{tab:kstar}
\begin{tabular}{lccc}
\toprule
$p$ & $\nu=3$ ($s=3/2$) & $\nu=4$ ($s=2$) & $\nu=5$ ($s=5/2$)\\
\midrule
59  & 5 & 4 & 3\\
71  & 5 & 4 & 3\\
113 & 7 & 4 & 4\\
\bottomrule
\end{tabular}
\end{table}

\paragraph{Scope.} Three hypotheses are load-bearing. (i) \emph{Full table}: Lemma~\ref{lem:junk}
averages over every input with a given answer; with a training fraction below one the training-set
margin need not obey the bound. (ii) \emph{Interference-freeness is a hypothesis on the table, not a
fact about trained networks}: networks trained with plain weight decay carry interference at the
sign-flipped modes, and \S\ref{sec:experiments} measures how much and how it moves with $q$.
(iii) $\Kstar$ is the support size of the free-share programme; the equal-share ladder is an
approximation, and where adjacent rungs are within a few percent the honest statement is a band of width
one. Whether any network \emph{outside} the class can exceed the class value is open; Appendix~\ref{app:conjecture}
states the conjecture and what is known under it.
